\section*{Introducción}
\addcontentsline{toc}{section}{Introducción}

El triage es un procedimiento de clasificación utilizado en los servicios de
urgencias para establecer el nivel de prioridad de atención de cada paciente
de acuerdo con su condición clínica. Su finalidad es permitir que los casos
de mayor gravedad sean atendidos antes, con independencia del orden de
llegada.

En la guardia pediátrica analizada los pacientes son clasificados en cuatro
categorías de prioridad --- rojo, naranja, amarillo y verde ---, ordenadas
desde la mayor hasta la menor urgencia. Esa clasificación condiciona el orden
en que los pacientes acceden a los recursos médicos disponibles.

El objetivo del trabajo es desarrollar un modelo de simulación que permita
representar el flujo de atención y analizar el impacto que las decisiones
tomadas durante el triage generan sobre los tiempos de espera, la permanencia
en el sistema y la utilización de los recursos. Asimismo, se busca emplear el
modelo para comparar el funcionamiento actual con un escenario alternativo y
evaluar posibles mejoras mediante técnicas de análisis estadístico.

El informe se organiza según los diez pasos de un estudio de simulación
vistos en la materia, que constituyen las secciones numeradas del documento.

